\documentclass[11pt,a4paper]{article}
\usepackage[paperwidth=210mm,paperheight=1100mm,margin=1.5cm]{geometry}
\usepackage{graphicx}
\usepackage{float}
\usepackage{caption}
\usepackage{subcaption}
\usepackage{amsmath}
\usepackage[hidelinks]{hyperref}
\usepackage{xcolor}

% No page numbers
\pagestyle{empty}
\makeatletter
\let\ps@plain\ps@empty
\makeatother

\setlength{\parindent}{0pt}
\setlength{\parskip}{0.5em}
\captionsetup{font=small,labelfont=bf}

% Adjust this in one place if the plots live somewhere else.
\newcommand{\plotdir}{../results}

\newcommand{\resultpair}[5]{%
\begin{figure}[H]
    \centering
    \begin{subfigure}[t]{0.48\linewidth}
        \centering
        \includegraphics[width=\linewidth]{\plotdir/#1}
        \caption{#2}
    \end{subfigure}
    \hfill
    \begin{subfigure}[t]{0.48\linewidth}
        \centering
        \includegraphics[width=\linewidth]{\plotdir/#3}
        \caption{#4}
    \end{subfigure}
    \caption{#5}
\end{figure}
}

\title{GNN Exercise: Neutrino Interaction Position Reconstruction}
\author{}
\date{}

\begin{document}

\maketitle

This is a quick summary of the final plots from the GNN model.

The model uses three DynamicEdgeConv layers, then global mean pooling, then a small MLP to predict the interaction point in the $(x,y)$ plane.

The plotting script only saves residual plots for $x$, so that is all included here.

\section*{Training Curve}

\begin{figure}[H]
    \centering
    \includegraphics[width=0.72\linewidth]{\plotdir/loss_curve.png}
    \caption{Training and validation MSE over epochs. Just a quick check that training behaves sensibly.}
    \label{fig:loss-curve}
\end{figure}


\section*{1. Overall Position Reconstruction}

\begin{figure}[H]
    \centering
    \includegraphics[width=0.72\linewidth]{\plotdir/test_predictions.png}
    \caption{True and predicted positions in the $(x,y)$ plane for the test set. This gives the overall picture.}
    \label{fig:test-predictions}
\end{figure}


\section*{2. Residual Diagnostics}

\resultpair
    {x_bias_vs_radius.png}
    {$x$ residual vs true radius.}
    {x_residual_vs_true_x.png}
    {$x$ residual vs true $x$.}
    {Two simple $x$ residual checks. The left plot looks for radial bias, and the right one looks for any clear left--right shift.}


\section*{3. Coordinate-wise Prediction Quality}

\resultpair
    {pred_vs_true_x.png}
    {Predicted vs true $x$. The dashed red line is perfect prediction.}
    {pred_vs_true_y.png}
    {Predicted vs true $y$. The dashed red line is perfect prediction.}
    {Predicted versus true values for each coordinate. Closer to the diagonal is better.}

\end{document}
